% ============================================================
%  VALIDATION CHAPTER
% ============================================================
\chapter{Validation}
\label{ch:validation}


% ============================================================
%  N.0  INTRODUCTION
% ============================================================
\section{Introduction}
\label{sec:val-intro}

The purpose of this chapter is to demonstrate that the analytical
framework developed in this thesis---FactorFlow extended with multi-layer
fusion support---accurately models the complex behavior of layer fusion
and hardware execution on a real accelerator design.
Establishing trust in the model is a prerequisite for the design-space
exploration performed in the subsequent chapters: if the analytical
predictions cannot be validated against a known hardware baseline, any
architectural insight drawn from them would be of limited value.

The validation is structured as follows.
Section~\ref{sec:val-methodology} summarizes the modeling methodology
and the role of Accelergy as validated energy-estimation back-end.
Section~\ref{sec:val-setup} describes the validation setup: the choice
of hardware architecture, how its characteristics were extracted from
the published paper, how the mapping and dataflow were abstracted into
the constraint model, and which workload was selected for comparison.
Section~\ref{sec:val-results} presents the validation results and
discusses the agreement between the model and the reference design.


% ============================================================
%  N.1  METHODOLOGY
% ============================================================
\section{Methodology}
\label{sec:val-methodology}

As described in Chapter~\ref{ch:proposed-solution}, the FactorFlow
framework was extended in Python to support multi-layer (fused)
computation.
The key extensions include:
\begin{enumerate}
  \item Generalized loop-nest representation that unifies all fused
        layers into a single loop nest with per-layer dimensions
        ($C_i$, $R_i$, $S_i$, $Z_i$, $Y_i$, $X_i$) and inter-layer
        coupling constraints.
  \item Memory-level bypass rules that distinguish between input, output,
        weight, intermediate-input, and intermediate-output operands,
        enabling precise modeling of on-chip data reuse across fused
        layers.
  \item Cost model that computes per-level memory operations (MOPs),
        energy, latency (including bandwidth-limited stall cycles), and
        energy--delay product (EDP) for the entire fused pipeline.
\end{enumerate}

Because this work is a modeling study without access to fabricated
silicon, energy per access at each level of the memory hierarchy is
estimated using Accelergy~\cite{accelergy}, the validated plug-in-based
energy-estimation framework developed by MIT.
Accelergy dispatches each component query to the most appropriate
estimator:
\begin{itemize}
  \item \textbf{CactiDRAM}~\cite{cacti} for off-chip DRAM (LPDDR4,
        64-bit bus);
  \item \textbf{CactiSRAM}~\cite{cacti} for on-chip SRAMs (Feature
        Memory, Weight Memory, GlobalBuffer);
  \item \textbf{Library plug-in} (based on Aladdin~\cite{aladdin}
        reference data with CMOS scaling~\cite{stillmaker2017scaling})
        for register files, multipliers, and adders;
  \item \textbf{Neurosim}~\cite{neurosim} for integer adders at
        process nodes not covered by Aladdin.
\end{itemize}
This energy-estimation pipeline is the same one used by
Timeloop~\cite{timeloop} and has been validated against post-layout
measurements and commercial SRAM compilers.


% ============================================================
%  N.2  VALIDATION SETUP
% ============================================================
\section{Validation Setup}
\label{sec:val-setup}


% ------------------------------------------------------------
\subsection{Rationale}
\label{sec:val-rationale}

The strongest form of validation for an analytical framework is
comparison against a real hardware implementation, rather than against
another high-level model.
A hardware-level comparison tests not only the mathematical formulation
of the cost model but also whether the architectural abstraction---memory
sizes, bandwidth, bypass rules, spatial mapping---faithfully captures the
behavior of the physical design.

For this reason, the validation target selected for this work is
\textbf{DepFiN}~\cite{goetschalckx2023depfin}, a depth-first CNN
accelerator fabricated in 22\,nm CMOS.
DepFiN is a particularly suitable reference for several reasons.

First, DepFiN \emph{natively supports layer fusion} as its primary
execution mode.
The architecture was designed from the ground up for depth-first
(fused) inference, with a split on-chip memory hierarchy (Feature
Memory and Weight Memory) and a fixed dataflow that pipelines
activations across up to ten consecutive convolutional layers without
intermediate DRAM accesses.
This makes it an ideal test case for the fusion extensions developed
in this thesis.

Second, DepFiN has been adopted as an architectural benchmark by
other analytical frameworks in the literature: DeFiNeS~\cite{defines}
and LoopTree~\cite{looptree} both model DepFiN as a representative
example of a depth-first fused accelerator, which confirms its
relevance and provides additional points of comparison.

Third, the DepFiN paper provides enough micro-architectural detail---PE
array dimensions, memory sizes, wordline widths, spatial mapping
strategy, and per-workload latency measurements---to construct a
faithful model, as described in the following sections.


% ------------------------------------------------------------
\subsection{Architecture Extraction}
\label{sec:val-arch-extraction}

\subsubsection{Parameters Taken from the Paper}

The key architectural parameters extracted from the DepFiN
paper~\cite{goetschalckx2023depfin} are summarized in
Table~\ref{tab:depfin-extracted-params}.

\begin{table}[htbp]
  \centering
  \caption{Architectural parameters extracted from the DepFiN
           paper~\cite{goetschalckx2023depfin}.}
  \label{tab:depfin-extracted-params}
  \small
  \begin{tabular}{l l l}
    \toprule
    \textbf{Parameter} & \textbf{Value} & \textbf{Source} \\
    \midrule
    Technology node         & 22\,nm                 & Paper Section~III \\
    Clock frequency         & $\sim$930\,MHz         & Paper Section~V \\
    PE array                & $16 \times 128 = 2048$ & Paper Figure~6 \\
    Feature Memory (FMEM)   & 1056\,KB (2 banks)     & Paper Section~IV \\
    FMEM wordline           & 1056\,bit              & Paper Section~IV \\
    Weight Memory (WMEM)    & 524\,KB (4 banks)      & Paper Section~IV \\
    WMEM wordline           & 128\,bit               & Paper Section~IV \\
    DRAM type               & LPDDR4                 & Paper Section~V \\
    DRAM I/O bandwidth      & 96\,bit = 12\,B/cycle  & Paper Figure~6 \\
    Operand precision       & 8\,bit                 & Paper Section~III \\
    Accumulator precision   & 32\,bit                & Paper Section~III \\
    Spatial mapping (cols)  & Output width ($Q$ / $X_i$) & Paper Section~IV \\
    Spatial mapping (rows)  & Output channels ($Z_i$)    & Paper Section~IV \\
    \bottomrule
  \end{tabular}
\end{table}

The 128 PE columns map onto the output spatial width dimension ($Q$ for
the final layer, $X_i$ for intermediate layers), while the 16 PE rows
map onto the output channel dimension ($Z_i$).
Because each layer produces 32 output channels and only 16 rows are
available, each tile requires two passes over the channel dimension
($Z_i/16 = 2$), which is reflected in the Weight Memory factor
constraints (\texttt{Z\_i: 2} for all intermediate layers).


\subsubsection{Assumptions and Limitations}
\label{sec:val-assumptions}

Not all parameters required for a complete analytical model are
available from the paper. The following assumptions were made:

\begin{enumerate}
  \item \textbf{DRAM bandwidth.}
        The paper reports an aggregate external I/O interface of
        96\,bits (12\,B/cycle) but does not provide a cycle-accurate
        DRAM bandwidth model; the reported latency values include
        unspecified stall effects from external memory.
        Precise DRAM bandwidth modeling would require knowledge of
        the LPDDR4 controller implementation, burst scheduling, and
        read/write turnaround penalties, none of which are disclosed.
        %
        For this reason, the validation focuses on the
        \emph{compute-bound} (or \emph{ideal}) latency---which depends
        solely on the number of MAC operations and the PE
        utilization---and on the on-chip memory hierarchy behavior,
        both of which are independent of the external memory bandwidth.

  \item \textbf{Feature Memory bandwidth.}
        We set the FMEM read bandwidth to 132\,B/cycle and write
        bandwidth to 128\,B/cycle, derived from the 1056-bit wordline
        and two-bank structure (1056\,bit / 8\,bit $\times$ 1 = 132\,B
        per read access).

  \item \textbf{Weight Memory bandwidth.}
        We set the WMEM bandwidth to 16\,B/cycle for both reads and
        writes, derived from the 128-bit wordline
        (128\,bit / 8\,bit = 16\,B per access).

  \item \textbf{Energy values.}
        Per-access energy values are estimated by Accelergy at 22\,nm
        rather than measured from the fabricated chip.
        The paper does not report per-component energy breakdowns, so
        direct comparison of energy is not possible; the validation
        focuses on latency and data-movement correctness.
\end{enumerate}


\subsubsection{Mapping Abstraction}
\label{sec:val-mapping}

Abstracting DepFiN's hardwired, custom dataflow into the generic
constraint model of FactorFlow was one of the most challenging aspects
of this work.
DepFiN is not a programmable accelerator with a flexible mapping
engine; it implements a \emph{fixed} depth-first dataflow in hardware,
where the data movement pattern is determined by the physical wiring
between the PE array, FMEM, and WMEM.

In FactorFlow, this fixed dataflow is expressed through
\textbf{dataflow constraints} and \textbf{factor constraints} at each
memory level.
The translation follows these principles:

\begin{enumerate}
  \item \textbf{DRAM level.}
        All spatial dimensions ($P$, $Q$, and all $Y_i$, $X_i$) are
        tiled at DRAM, reflecting the fact that the accelerator
        processes one spatial tile at a time before moving to the next.
        The output width is divided into tiles of 128 pixels (matching
        the 128 PE columns), yielding $Q = 1280/128 = 10$ tiles.
        All 720 rows are iterated at this level ($P = 720$, $Y_i = 720$).
        Intermediate operands (\texttt{int\_in}, \texttt{int\_out}) bypass
        DRAM entirely, enforcing the depth-first fusion invariant that
        no intermediate data touches external memory.

  \item \textbf{Feature Memory.}
        FMEM stores all non-weight operands: input activations, output
        activations, and intermediate activations between fused layers.
        Weights bypass FMEM.
        No additional tiling factors are constrained at this level,
        allowing the mapper to determine the optimal reuse order.

  \item \textbf{Weight Memory.}
        WMEM stores all weights for all 11 layers (10 fused +~1 output)
        simultaneously.
        All weight-related dimensions ($C_i$, $R_i$, $S_i$, $Z_i$)
        are constrained to their full values at this level, reflecting
        the fact that the complete filter set is pre-loaded before
        inference begins.
        All activation operands bypass WMEM.

  \item \textbf{Spatial fanout levels.}
        Each of the 11 layers receives a dedicated pair of spatial
        fanout levels: \texttt{SACols\_$i$} (mesh\,=\,128) distributes
        the output width, and \texttt{SARows\_$i$} (mesh\,=\,16)
        distributes the output channels.
        The factor constraints ensure that exactly 128 positions and
        16 channels are processed in parallel per tile.

  \item \textbf{Accumulation register.}
        A single \texttt{AccumulationIntermediateOutputRegister}
        sits below all spatial levels and above the compute units.
        All spatial and output dimensions are constrained to factor~1
        at this level, meaning it stores exactly one partial sum per
        PE at a time.
        This register handles both output accumulation (across input
        channels) and intermediate-output storage (between fused
        layers).
\end{enumerate}

This constraint-based encoding faithfully reproduces DepFiN's
hardwired dataflow: the mapper has no freedom to reorder the
weight-related or spatial dimensions, and the only degrees of freedom
are the tiling and reuse ordering within FMEM---exactly as in the
physical design.


% ------------------------------------------------------------
\subsection{Workload Selection}
\label{sec:val-workload}

The DepFiN paper~\cite{goetschalckx2023depfin} evaluates several
workloads in its measurement results.
For validation, we selected \textbf{Workload~2}, a 10-layer fusion
pipeline operating on 720$\times$1280 HD-resolution
frames.
Its layer structure is shown in Table~\ref{tab:val-workload}.

\begin{table}[htbp]
  \centering
  \caption{Workload~2 from the DepFiN paper: 10-layer fused pipeline
           at 720$\times$1280 resolution.
           Layers~1--9 are identical 3$\times$3 convolutions with
           32 input and 32 output channels.}
  \label{tab:val-workload}
  \small
  \begin{tabular}{c c c c c c}
    \toprule
    \textbf{Layer} & \textbf{Resolution} & \textbf{Channels (in$\to$out)}
      & \textbf{Kernel} & \textbf{MACs} \\
    \midrule
    L0  & 720$\times$1280 & 3$\to$32    & 3$\times$3 &   796,262,400 \\
    L1  & 720$\times$1280 & 32$\to$32   & 3$\times$3 & 8,493,465,600 \\
    L2  & 720$\times$1280 & 32$\to$32   & 3$\times$3 & 8,493,465,600 \\
    L3  & 720$\times$1280 & 32$\to$32   & 3$\times$3 & 8,493,465,600 \\
    L4  & 720$\times$1280 & 32$\to$32   & 3$\times$3 & 8,493,465,600 \\
    L5  & 720$\times$1280 & 32$\to$32   & 3$\times$3 & 8,493,465,600 \\
    L6  & 720$\times$1280 & 32$\to$32   & 3$\times$3 & 8,493,465,600 \\
    L7  & 720$\times$1280 & 32$\to$32   & 3$\times$3 & 8,493,465,600 \\
    L8  & 720$\times$1280 & 32$\to$32   & 3$\times$3 & 8,493,465,600 \\
    L9  & 720$\times$1280 & 32$\to$32   & 3$\times$3 & 8,493,465,600 \\
    L10 & 720$\times$1280 & 32$\to$16   & 1$\times$1 &   471,859,200 \\
    \midrule
    \multicolumn{4}{r}{\textbf{Total}} & \textbf{77,709,312,000} \\
    \bottomrule
  \end{tabular}
\end{table}

This workload was chosen for three reasons:

\begin{enumerate}
  \item \textbf{Minimal external effects.}
        Of the workloads evaluated in the DepFiN paper, Workload~2 has
        the fewest dependencies on implementation-specific optimizations
        that are not captured by the analytical model, such as power
        gating of unused PEs, dynamic clock scaling, or partial-array
        configurations described in the ``Measurements and Metrics''
        section of~\cite{goetschalckx2023depfin}.
        Workloads that rely on these techniques would introduce
        discrepancies unrelated to the modeling accuracy of the
        dataflow and cost model.

  \item \textbf{Full PE utilization.}
        With 32 output channels distributed across 16 PE rows
        ($Z_i/\text{rows} = 32/16 = 2$) and 1280 output width pixels
        distributed across 128 PE columns
        ($Q/\text{cols} = 1280/128 = 10$), every PE is utilized at
        every cycle during computation.
        The spatial utilization is exactly 1.0 (100\,\%), as confirmed
        by the model output.
        This eliminates PE idle time as a source of modeling error.

  \item \textbf{Homogeneous structure.}
        Nine of the eleven layers (L1--L9) are identical
        3$\times$3 convolutions with 32 input and 32 output channels.
        This regularity simplifies the analysis and makes discrepancies
        easier to identify and attribute to specific modeling choices.
\end{enumerate}


% ============================================================
%  N.3  VALIDATION RESULTS
% ============================================================
\section{Validation Results}
\label{sec:val-results}

The validation run was executed as:
\begin{center}
\texttt{python3 main\_cli.py architectures/multi\_layer\_arch.py architectures/10layer\_comp.py}
\end{center}

\noindent Table~\ref{tab:val-results} reports the key results.

\begin{table}[htbp]
  \centering
  \caption{Validation results: FactorFlow model vs.\ DepFiN ideal
           (compute-bound) reference for Workload~2.}
  \label{tab:val-results}
  \small
  \begin{tabular}{l r r}
    \toprule
    \textbf{Metric}
      & \textbf{DepFiN (ideal)}
      & \textbf{FactorFlow model} \\
    \midrule
    Total MAC operations
      & \multicolumn{2}{c}{77,709,312,000} \\
    PE count
      & \multicolumn{2}{c}{2,048 ($16 \times 128$)} \\
    Spatial utilization
      & 100\,\%
      & 100\,\% \\
    \addlinespace
    Ideal latency (cycles)
      & $\sim$37.94\,M
      & 37,944,000 \\
    \addlinespace
    FMEM stall cycles
      & ---
      & 24 \\
    WMEM stall cycles
      & ---
      & 0 \\
    \addlinespace
    DRAM intermediate reads
      & 0
      & 0 \\
    DRAM intermediate writes
      & 0
      & 0 \\
    \bottomrule
  \end{tabular}
\end{table}


\subsection{Ideal Latency}
\label{sec:val-ideal-latency}

The ideal compute-bound latency is obtained by dividing the total
number of MAC operations by the number of PEs:
%
\begin{equation}
  \label{eq:ideal-latency}
  L_{\text{ideal}} = \frac{\text{Total MACs}}{\text{PE count}}
  = \frac{77{,}709{,}312{,}000}{2{,}048}
  = 37{,}944{,}000 \;\text{cycles}
  \approx 37.94\,\text{M cycles}
\end{equation}
%
This value represents the theoretical minimum latency assuming
100\,\% PE utilization and no memory stalls---i.e., all data is
available at the compute units without any bandwidth bottleneck.

The FactorFlow model produces exactly this value as the ideal
per-layer latency summed across all layers.
Specifically, for each of the nine identical 3$\times$3 convolutional
layers (L1--L9), the model computes:
%
\begin{equation}
  L_{\text{per-layer}} = \frac{720 \times 1280 \times 32 \times 32 \times 3 \times 3}{2{,}048}
  = 4{,}147{,}200 \;\text{cycles}
\end{equation}
%
while Layer~0 (3 input channels) contributes 388,800 cycles and
Layer~10 (1$\times$1 convolution, 16 output channels) contributes
230,400 cycles, for a total of $9 \times 4{,}147{,}200 + 388{,}800 +
230{,}400 = 37{,}944{,}000$ cycles.
This exact agreement with the reference value, computed directly from
the workload parameters and PE count, confirms that the unified
loop-nest formulation correctly accounts for all MAC operations across
the fused layers, with no under- or over-counting.


\subsection{Data Residency and Intermediate Buffers}
\label{sec:val-data-residency}

A central claim of the DepFiN design is that, in fused mode,
\textbf{intermediate activations between consecutive layers never
touch DRAM}.
The model confirms this property: the total DRAM intermediate reads
and writes are both exactly zero (Table~\ref{tab:val-results}).

This result is a direct consequence of the bypass constraints in the
architectural model:
\begin{itemize}
  \item DRAM bypasses \texttt{int\_in} and \texttt{int\_out}, so
        intermediate operands cannot be stored in or fetched from
        external memory.
  \item FMEM stores all intermediate activations, as reflected by the
        per-layer intermediate-output read and write counts of
        29,491,200 each for layers L1--L9 (and 825,753,600 for L0's
        intermediate writes consumed by successive layers).
\end{itemize}

The model further shows that the 1056\,KB Feature Memory provides
sufficient capacity to hold all intermediate data for the 10-layer
pipeline without eviction.
The FMEM footprint remains fully cached: the working set of
intermediate activations (one tile's output from one layer = one
tile's input for the next layer) fits within the 1056\,KB budget,
so no intermediate data is ever spilled to or re-fetched from DRAM.
This validates the architectural design point of DepFiN: the FMEM size
is indeed sufficient for full depth-first fusion of the target
workload, and no intermediate recomputation is required.


\subsection{On-Chip Memory Stalls}
\label{sec:val-stalls}

A key validation result concerns the stall behavior of the
\emph{on-chip} memory hierarchy, which can be fully characterized
because the FMEM and WMEM bandwidths are derived directly from the
known wordline widths (Section~\ref{sec:val-assumptions}).

\paragraph{Weight Memory.}
The Weight Memory introduces \textbf{zero stall cycles} across all
eleven layers.
The 16\,B/cycle bandwidth (128-bit wordline) is sufficient to deliver
weights to the PE array without ever becoming the bottleneck.

\paragraph{Feature Memory.}
The Feature Memory introduces a total of only \textbf{24 stall cycles}
across the entire 37.94\,M-cycle execution---a negligible overhead of
$6.3 \times 10^{-5}\,\%$.
Crucially, the per-layer breakdown reveals that these 24 stall cycles
arise \textbf{exclusively from the two boundary layers}:

\begin{itemize}
  \item \textbf{Layer~0} (input layer, $C_0 = 3$): contributes 19
        stall cycles.
        This layer has only 3 input channels instead of 32, which
        alters the ratio between the data volume to be read from FMEM
        and the compute time, causing a minor bandwidth mismatch
        during the first few tiles.
  \item \textbf{Layer~10} (output layer, 1$\times$1 convolution,
        $Z_{10} = 16$): contributes 5 stall cycles.
        This layer has a reduced kernel size and output channel count
        compared to the regime layers, creating a similarly marginal
        imbalance.
  \item \textbf{Layers~1--9} (regime layers): \textbf{zero stall
        cycles each}.
        All nine identical 3$\times$3 convolutions with 32 input and
        32 output channels execute without any FMEM bandwidth
        bottleneck.
\end{itemize}

This result has two important implications.
First, it confirms that the FMEM read bandwidth (132\,B/cycle) and
write bandwidth (128\,B/cycle) are well-dimensioned for the target
dataflow: at regime---where the pipeline spends the vast majority of
its execution time---the on-chip data delivery perfectly sustains the
PE throughput with zero stalls.
Second, the negligible stalls at the boundary layers are a
\emph{structural} artefact of the non-uniform layer shapes (3 input
channels for L0, 1$\times$1 kernel for L10) rather than a systematic
bandwidth limitation.
In a longer pipeline with more regime layers, these boundary effects
would become even more insignificant relative to the total execution
time.


\subsection{DRAM Bandwidth}
\label{sec:val-dram-bw}

As discussed in Section~\ref{sec:val-assumptions}, the DepFiN paper
does not provide enough information to construct a cycle-accurate
model of the DRAM bandwidth.
The reported aggregate I/O interface of 12\,B/cycle (96\,bits) does
not account for LPDDR4 controller overhead, read/write turnaround
penalties, or burst scheduling effects.
For this reason, we do not attempt to validate the absolute total
latency (which would include DRAM-induced stalls) and instead focus
the comparison on the compute-bound latency and the on-chip memory
behavior---quantities that are fully determined by the architectural
parameters known from the paper.

The model nonetheless confirms that, regardless of the DRAM bandwidth
assumption, the depth-first fusion guarantee holds: all intermediate
data remains on-chip, and the only DRAM traffic consists of the
unavoidable input reads (2,764,800 bytes for the input image) and
output writes (14,745,600 bytes for the final output), plus the
one-time weight loads (84,320 bytes).


\subsection{Summary}

The validation demonstrates that the FactorFlow model with layer-fusion
extensions faithfully reproduces the key properties of the DepFiN
accelerator:
\begin{enumerate}
  \item The \textbf{ideal latency} of 37.94\,M cycles matches exactly
        the compute-bound reference
        ($\text{Total MACs}/\text{PEs}$), confirming correct
        MAC-operation accounting across 11 fused layers.
  \item The \textbf{DRAM bypass} of intermediate operands is correctly
        enforced, with zero intermediate DRAM traffic, reproducing
        DepFiN's depth-first fusion guarantee.
  \item The \textbf{FMEM capacity} is verified to be sufficient for
        the entire fused pipeline, with no intermediate data
        spills---confirming that the 1056\,KB Feature Memory provides
        fully-cached intermediate storage.
  \item The \textbf{on-chip memory stalls are negligible}: the Feature
        Memory introduces only 24 stall cycles out of 37.94\,M, and
        these arise exclusively from the boundary layers (L0 and L10).
        At regime (layers L1--L9), the FMEM bandwidth sustains the PE
        throughput with \textbf{zero stall cycles}, validating that the
        on-chip data delivery is perfectly matched to the compute
        demand.
  \item The \textbf{Weight Memory} introduces zero stall cycles,
        confirming that the 16\,B/cycle bandwidth is sufficient for
        weight delivery across all layers.
\end{enumerate}

These results establish confidence in the analytical framework for the
design-space exploration presented in the following chapters.
The model correctly captures the fundamental trade-offs of layer
fusion---DRAM traffic elimination, on-chip memory pressure, and
bandwidth utilization---which are the quantities that drive the
energy--latency comparisons in the case studies.
